% ==============================================================================
% EM tenure with duration contamination (rho model)
% Created: 2026-04-29
%
% This document extends "EM tenure.tex" with a duration contamination parameter
% rho that separates tenure reporting error from employment misclassification.
% The original document is NOT modified; this is a self-contained companion.
% ==============================================================================
\documentclass[12pt]{article}
\usepackage[margin=2cm]{geometry}
\setlength{\parindent}{1em}
\setlength{\parskip}{1em}

\usepackage{amsmath}
\usepackage{booktabs,dcolumn,setspace, float}

\begin{document}

\section{Duration Contamination Extension of the EM Framework}

This document presents an extension of the EM framework in ``EM tenure.tex'' that introduces a duration contamination probability~\(\rho\). The extension addresses the empirical finding that approximately 37\% of tenure increments in the QLFS are inconsistent with the clock model (non-standard increments), even though the individuals are correctly classified as employed. The base model attributes these to state misclassification, inflating~\(\pi\) from~\(\sim\!3\%\) to~\(\sim\!25\%\). The~\(\rho\) model provides a separate error channel for ``correct state, wrong duration.''

All notation, conventions, and base-model equations follow ``EM tenure.tex'' exactly. We describe only the changes and additions here.

\subsection{Motivation}

In the South African QLFS, 63\% of employed\(\to\)employed tenure increments are exactly 0.25~years (to machine precision), consistent with the clock model \(\Delta^g_t = 0\) and \(\sigma^2_g = 0\). The remaining 37\% include 15\% negative increments (median \(-21\)~months, far exceeding any plausible within-quarter job change) and 22\% positive but non-standard increments. These are the signature of \emph{tenure reporting error}: the same person at the same job reports tenure differently across waves due to rounding, recall bias, or respondent substitution.

The base model has one error channel (\(\pi\)) that handles employment state misclassification. It has no mechanism for ``correct state, wrong duration.'' When \(\sigma^2_g \to 0\) (as the MLE correctly finds), all non-standard increments must be explained through misclassification histories, inflating~\(\pi\).

\subsection{Extended observation model}

\paragraph{Duration contamination.}
We introduce a per-wave latent indicator \(m_{it} \in \{0,1\}\) with
\[
\Pr(m_{it} = 1) = \rho, \qquad \Pr(m_{it} = 0) = 1 - \rho.
\]
When \(m_{it} = 0\) (``clock-consistent''), the duration emission follows the base model. When \(m_{it} = 1\) (``contaminated''), the reported duration is an uninformative draw from the population marginal, independent of the individual's trajectory.

\paragraph{Modified per-wave emission.}
The per-wave contribution to the individual likelihood becomes:
\begin{equation}
\ell^*_t(y_{it} \mid h, \Theta) = q(s_{it} \mid h_t, \pi) \times \Big[(1 - \rho)\,\ell^{\mathrm{dur}}_t(y_{it} \mid h, s_{it}, \Theta) + \rho\,f^{\mathrm{dur}}_{\mathrm{pop}}(y_{it} \mid s_{it}, \Theta)\Big],
\label{eq:emission_rho}
\end{equation}
where:
\begin{itemize}
\item \(q(s_{it} \mid h_t, \pi)\) is the misclassification probability from the base model, unchanged;
\item \(\ell^{\mathrm{dur}}_t\) is the existing duration emission (Normal increment, EMG, interval-censored, etc.)\ --- exactly the base model's duration component;
\item \(f^{\mathrm{dur}}_{\mathrm{pop}}\) is the population marginal for the observed duration given the observed state:
\begin{align}
f^{\mathrm{dur}}_{\mathrm{pop}}(y_{it} \mid s_{it} = 1, \Theta) &= f_{\mathrm{EMG}}(g_{it};\, \lambda_g,\, \sigma^2_g), \label{eq:fpop_g} \\
f^{\mathrm{dur}}_{\mathrm{pop}}(y_{it} \mid s_{it} = 0, \Theta) &= P(D \in [a_{k_{it}}, b_{k_{it}}) \mid \lambda_d) \quad \text{(interval-censored Exp)}. \label{eq:fpop_d}
\end{align}
\end{itemize}

\paragraph{Key separation.} Misclassification (\(\pi\)) affects the \emph{state observation}: \(P(s_t \neq h_t) = \pi\). Duration contamination (\(\rho\)) affects the \emph{duration observation}: with probability~\(\rho\), the reported duration is uninformative about the trajectory. The two error channels operate on different margins of the data and do not interact in the complete-data log-likelihood.

\subsection{Observed-data likelihood}

The observed-data log-likelihood under the extended model is:
\begin{equation}
\ell(\Theta) = \sum_{i=1}^N w_i \log\bigg[\sum_{h \in \mathcal{H}} p(H_i = h \mid \alpha, \theta_0, \theta_1) \prod_{t=1}^3 \ell^*_t(y_{it} \mid h, \Theta)\bigg],
\label{eq:obs_ll_rho}
\end{equation}
with parameter vector \(\Theta = (\alpha, \theta_0, \theta_1, \pi, \rho, \sigma^2_g, \lambda_g, \lambda_d)\).

\subsection{Complete-data formulation}

Introduce both the history indicators \(z_{ih} = \mathbf{1}\{H_i = h\}\) and the contamination indicators \(m_{it} \in \{0,1\}\). The complete-data log-likelihood is:
\begin{align}
\log L_c(\Theta) &= \sum_{i=1}^N \sum_{h \in \mathcal{H}} w_i z_{ih} \bigg\{
\underbrace{\log p(h \mid \alpha, \theta_0, \theta_1)}_{\text{Markov prior}} \nonumber\\
&\quad + \underbrace{\sum_{t=1}^3 \log q(s_{it} \mid h_t, \pi)}_{\text{state misclassification}} \nonumber\\
&\quad + \underbrace{\sum_{t=1}^3 \Big[(1 - m_{it}) \log \ell^{\mathrm{dur}}_t + m_{it} \log f^{\mathrm{dur}}_{\mathrm{pop}}\Big]}_{\text{duration emission (mixture)}} \nonumber\\
&\quad + \underbrace{\sum_{t=1}^3 \Big[m_{it} \log \rho + (1 - m_{it}) \log(1 - \rho)\Big]}_{\text{contamination prior}}
\bigg\}.
\label{eq:complete_ll_rho}
\end{align}

The four blocks in~\eqref{eq:complete_ll_rho} separate cleanly: the Markov prior depends on \((\alpha, \theta_0, \theta_1)\), misclassification on~\(\pi\), the duration mixture on \((\sigma^2_g, \lambda_g, \lambda_d)\), and the contamination prior on~\(\rho\).

\subsection{E-step}

\paragraph{Step 1: Responsibilities.}
Compute \(\gamma_{ih}\) as in the base model, but using the mixture emission~\eqref{eq:emission_rho}:
\begin{equation}
\gamma_{ih} \propto p(h \mid \Theta^{\mathrm{old}}) \prod_{t=1}^3 \ell^*_t(y_{it} \mid h, \Theta^{\mathrm{old}}).
\label{eq:resp_rho}
\end{equation}

\paragraph{Step 2: Posterior contamination probabilities.}
For each observation~\(i\), history~\(h\), and wave~\(t\), compute the posterior probability that the duration is contaminated:
\begin{equation}
\omega_{ith} = \frac{\rho^{\mathrm{old}} \cdot f^{\mathrm{dur}}_{\mathrm{pop}}(y_{it} \mid s_{it}, \Theta^{\mathrm{old}})}{(1 - \rho^{\mathrm{old}}) \cdot \ell^{\mathrm{dur}}_t(y_{it} \mid h, \Theta^{\mathrm{old}}) + \rho^{\mathrm{old}} \cdot f^{\mathrm{dur}}_{\mathrm{pop}}(y_{it} \mid s_{it}, \Theta^{\mathrm{old}})}.
\label{eq:omega}
\end{equation}
This is the standard Bayesian posterior for a two-component mixture. When \(\ell^{\mathrm{dur}}_t\) is large relative to \(f^{\mathrm{dur}}_{\mathrm{pop}}\) (e.g., \(\Delta g_t = 0.25\) exactly), \(\omega_{ith} \approx 0\). When \(\ell^{\mathrm{dur}}_t \approx 0\) (wildly inconsistent duration), \(\omega_{ith} \approx 1\).

\paragraph{Modified sufficient statistics.}
The Markov transition statistics \((C_1, C_0, D_1, D_0, T_{11}, T_{01})\) and the misclassification count~\(M\) are computed from \(\gamma_{ih}\) exactly as in the base model --- they do \emph{not} use~\(\omega_{ith}\).

The duration variance statistics are modified to use \((1 - \omega)\) weights:
\begin{align}
\tilde{S}_g &= \sum_{i,h} w_i \gamma_{ih} \sum_{t:\,\text{cont.}} (1 - \omega_{ith})\, (\Delta^g_{it})^2, \qquad
\tilde{N}_g = \sum_{i,h} w_i \gamma_{ih} \sum_{t:\,\text{cont.}} (1 - \omega_{ith}), \label{eq:suff_sg_rho} \\
\tilde{S}^{\mathrm{start}}_g &= \sum_{i,h} w_i \gamma_{ih} \sum_{t:\,\text{start}} (1 - \omega_{ith})\, (g_{it} - 0.25)^2, \qquad
\tilde{N}^{\mathrm{start}}_g = \sum_{i,h} w_i \gamma_{ih} \sum_{t:\,\text{start}} (1 - \omega_{ith}). \label{eq:suff_sg_start_rho}
\end{align}

The contamination sufficient statistic is:
\begin{equation}
\Omega = \sum_{i,h} w_i \gamma_{ih} \sum_{t=1}^3 \omega_{ith}.
\label{eq:suff_omega}
\end{equation}

\subsection{M-step}

\paragraph{Markov prior block} (\(\alpha, \theta_0, \theta_1\)): unchanged from the base model. Uses \(\gamma_{ih}\) only.
\begin{equation}
\hat{\alpha} = \frac{C_1}{C_0 + C_1}, \qquad \hat{\theta}_1 = \frac{T_{11}}{D_1}, \qquad \hat{\theta}_0 = \frac{T_{01}}{D_0}.
\end{equation}

\paragraph{Misclassification} (\(\pi\)): unchanged.
\begin{equation}
\hat{\pi} = \frac{M}{3W}, \qquad W = \sum_i w_i.
\end{equation}

\paragraph{Duration contamination} (\(\rho\)): closed-form.
\begin{equation}
\hat{\rho} = \frac{\Omega}{3W},
\label{eq:mstep_rho}
\end{equation}
where \(\Omega\) is from~\eqref{eq:suff_omega}. This is the weighted average posterior contamination probability.

\paragraph{Measurement variance for tenure} (\(\sigma^2_g\)): uses \((1 - \omega)\)-weighted statistics.
\begin{equation}
\hat{\sigma}^2_g = \frac{\tilde{S}_g + 2\,\tilde{S}^{\mathrm{start}}_g}{2(\tilde{N}_g + \tilde{N}^{\mathrm{start}}_g)}.
\label{eq:mstep_sigma_g_rho}
\end{equation}

\paragraph{Exponential rates} (\(\lambda_g, \lambda_d\)): The Brent root-finding for \(\lambda_g\) and \(\lambda_d\) proceeds exactly as in the base model. All EMG-type observations contribute --- both clock-consistent and contaminated use the same distributional form \(f_{\mathrm{EMG}}(x; \lambda_g, \sigma^2_g)\), so the FOC is unchanged. The EMG weights for each observation become \(w_i \gamma_{ih}\) (summing over both \(m_{it} = 0\) and \(m_{it} = 1\) contributions, since both involve the same EMG density).

\subsection{Identification}
\label{sec:identification_rho}

\paragraph{\(\pi\) vs \(\rho\): separate margins.}
\(\pi\) is identified from the \emph{state observation}: the misclassification probability \(q(s_t \mid h_t, \pi) = \pi^{\mathbf{1}\{s_t \neq h_t\}}\!(1 - \pi)^{\mathbf{1}\{s_t = h_t\}}\) does not involve any duration data. The M-step update \(\hat{\pi} = M / (3W)\) uses only the state disagreement count~\(M\), which depends on the observed states~\(s_t\) and the latent histories~\(h_t\) through \(\gamma_{ih}\) but does \emph{not} depend on~\(\omega_{ith}\).

\(\rho\) is identified from the \emph{duration observation}: the posterior \(\omega_{ith}\) compares the clock-consistent density \(\ell^{\mathrm{dur}}_t\) to the population marginal \(f^{\mathrm{dur}}_{\mathrm{pop}}\). When the increment is exactly 0.25~years, \(\ell^{\mathrm{dur}}_t\) is very large (near a point mass as \(\sigma^2_g \to 0\)) while \(f^{\mathrm{dur}}_{\mathrm{pop}}\) is moderate, so \(\omega_{ith} \approx 0\). When the increment deviates wildly, \(\ell^{\mathrm{dur}}_t \approx 0\) while \(f^{\mathrm{dur}}_{\mathrm{pop}}\) is reasonable, so \(\omega_{ith} \approx 1\).

\paragraph{How durations inform \(\theta_1\) and \(\pi\) through \(\gamma_{ih}\).}
Although the M-step updates for \(\theta\) and \(\pi\) do not use~\(\omega\) directly, the \emph{E-step responsibilities}~\(\gamma_{ih}\) depend on the duration data through the mixture emission~\eqref{eq:emission_rho}. This is the key channel by which durations inform transition and misclassification estimates:
\begin{itemize}
\item For an individual observed (1,\,0,\,1) with \(g_3 = 0.25\): tenure3 is consistent with a fresh start, so \(\gamma\) is pushed toward \(h = (1,0,1)\) (genuine transition, not misclassification).
\item For an individual observed (1,\,0,\,1) with \(g_3 = 5.5\): tenure3 is consistent with continuous employment, so \(\gamma\) is pushed toward \(h = (1,1,1)\) (state misclassification at wave~2).
\end{itemize}
The \(\rho\) model preserves this channel while discounting the duration signal for observations where the duration is clearly contaminated (high~\(\omega\)).

\subsection{Nesting and testing}

The base model is nested: setting \(\rho = 0\) in~\eqref{eq:emission_rho} recovers the base emission \(\ell^{\mathrm{dur}}_t\) exactly. The likelihood ratio test
\[
\mathrm{LR} = 2\big[\ell_{\rho} - \ell_{\text{base}}\big] \sim \chi^2_1 \quad \text{under } H_0: \rho = 0
\]
tests whether duration contamination is statistically significant. One degree of freedom because the only additional parameter is~\(\rho\).

\subsection{Summary of parameterisations}

\begin{center}
\begin{tabular}{llll}
\toprule
Block & Base model & \(\rho\)-extended model & Notes \\
\midrule
\(\alpha\) & Closed form & Same & Stationarity option. \\
\((\theta_1, \theta_0)\) & Closed form & Same & From prior only. \\
\(\pi\) & Closed form & Same & From state disagreements. \\
\(\rho\) & --- & Closed form \eqref{eq:mstep_rho} & New: from \(\omega_{ith}\). \\
\(\sigma^2_g\) & Closed form & \((1-\omega)\)-weighted \eqref{eq:mstep_sigma_g_rho} & Clock-consistent obs only. \\
\(\lambda_g\) & Brent & Same & All EMG obs contribute. \\
\(\lambda_d\) & Brent & Same & All interval obs contribute. \\
\bottomrule
\end{tabular}
\end{center}

The \(\rho\)-extended M-step retains the clean block structure of the base model, adding one closed-form update for~\(\rho\) and replacing the raw sufficient statistics for~\(\sigma^2_g\) with their \((1-\omega)\)-weighted counterparts. All other updates are unchanged.

\end{document}
